\documentclass[a4paper,12pt]{extarticle}
\usepackage[margin=1.5cm]{geometry}
\usepackage{amsmath, amssymb}
\usepackage{array}
\usepackage{multicol}
\usepackage{helvet}
\usepackage{multirow}
\usepackage{graphicx}
\renewcommand{\baselinestretch}{2}
\renewcommand{\familydefault}{\sfdefault}

\begin{document}

    \begin{tabular}{| c | c | c | c |}
        \hline
        & MATEMÁTICAS & 17/03/2026 & calificación \\
        \hline 

        \multirow{2}{*}{\includegraphics{../../img-Logos/iesoGalileo.jpg}} & Nombre y apellidos \hspace{120pt} & Ex. 1º ESO UD5:  & \\ 
        &&Álgebra& \\ 
        \hline
    \end{tabular}
    
    \begin{tabular}{| l | }
        \hline
        NOTAS: \\ 
            1. El examen tiene que ser limpio, ordenado y sin faltas de ortografía. \\
            
            2. El examen ha de realizarse en bolígrafo negro o azul, evitando tachones en la medida \\
            de lo posible. \\ 
            
            3. Deben aparecer todas las operaciones NO vale con indicar solo el resultado.\\
            
            4. Los problemas deben contener: Datos, planteamiento y resolución, respondiendo a lo \\
            que se pregunte, NO vale con indicar un número como solución del problema.\\

            5. Persona que se vea copiando o hablando con un compañero tiene un cero en el examen\\

            6. NO está permitido el uso de calculadoras \\
        \hline
    \end{tabular}

    \begin{enumerate}
        \item (1,5 puntos) Calcula dando el resultado lo más simplificado posible:
        \[3+\dfrac{2}{7} \cdot (1-\dfrac{1}{3})\] $\vspace{4cm}$ 
        \item (1 punto)  Alberto ha fallado 2 penaltis de 5 y Carlos ha marcado 4 de 7. ¿Quién tira mejor los penaltis?$\vspace{5cm}$ 
        \item (0,5 puntos) Completa las siguientes fracciones para que sean equivalentes.
        \[ \dfrac{7}{x} = \dfrac{14}{4} = \dfrac{y }{6}\] $\vspace{5cm}$ 
        \item (3 puntos) Realiza las siguientes operaciones
        \begin{enumerate}
            \item 19.4 + 28.78 - 30.791=
            \item $10.25 \cdot 8.7 = $
            \item 0.23 : 1.2 = $\vspace{5cm}$ 
        \end{enumerate}
        \item (1 punto) Francisco ha comprado tres bolígrafos y dos rotuladores. Un bolígrafo cuesta 0.45 € y un rotulador 1.20 € ¿Cuánto le devuelven si paga con un billete de 5 €?$\vspace{5cm}$ 
        \item (1 punto) Resuelve la siguiente Ecuación: 
        \[ 5x = 4 - 3x + 5 - x\]$\vspace{3cm}$ 
        \item (1 punto) Completa la tabla con la siguiente información: 
        
        \begin{itemize}
            \item Cristina tiene x años
            \item Alberto, su esposo, tiene 3 años más 
            \item Javier, su padre, le dobla la edad
            \item Marta, su madre, tiene 4 años menos que su padre
            \item Loli y Mar son sus hijas gemelas. Las tuvo con 26 años.
            \item Javi, el pequeño, tiene la mitad de años que las gemelas.
        \end{itemize} 
        \begin{tabular}{| c | c | c | c | c | c | c | }
                \hline
                & Cristina & Alberto & Javier & Marta & Loli y Mar &  Javi \\ \hline
                Edad & x & $ \hspace{2cm}$  & $ \hspace{2cm}$ & $ \hspace{2cm}$  &$ \hspace{2cm}$ &  $ \hspace{2cm}$\\ 
                &  &  &  &   & &  \\ 
                \hline
            \end{tabular}
        \item (1 punto) Luis tiene dos años más que su hermana y entre los dos suman 18 años. ¿Cuál es la edad de cada hermano? 

    \end{enumerate}
\end{document}